\chapter{Обзор предметной области}
Использование оптических технологий и элементов является большим прогрессом в наше время. Меньшая длина волны позволяет достичь большей частоты сигналов, в частности при передаче данных. Такая частота дает возможность передавать потоки информации в несколько терабит в секунду. Важными преимуществами ВОЛС являются такие факторы, как малое затухание сигналов, позволяющее, при использовании современных технологий, строить участки оптических систем в сто и более километров без ретрансляции, высокая помехозащищенность, связанная с малой восприимчивостью оптического волокна к электромагнитным помехам, и многие другие.

\begin{figure}[h!]
\includegraphics[width=0.5\textwidth]{img/refraction.pdf}
\caption{Преломление света}
\end{figure}

Волоконная оптика - это раздел оптики, изучающий процессы в оптоволокне. Оптическое волокно — нить из оптически прозрачного материала (стекло, пластик), используемая для переноса света внутри себя посредством полного внутреннего отражения. Для объяснения этого явления рассмотрим процессы происходящие при прохожении светом границы двух сред:

При переходе в другую среду свети отклоняется от направления своего распространения. Величина этого отклонения определяется законом Снеллиуса:
\begin{equation}
 n_1 \sin\theta_1 = n_2 \sin\theta_2
\end{equation}

Очевидно, что при некотором угле падения синус соответствуюшего угла преломленния окажется больше единицы, что невозможно. В этом случае будет наблюдаться эффект полного внутреннего отражения. Это возможно только в случае $n_1 > n_2$

\begin{figure}[h!]
\includegraphics[width=0.5\textwidth]{img/refraction_full_inner.pdf}
\caption{Полное внутреннее отражение}
\end{figure}

Планарный волновод - это структура из слоев с тремя разными показателями преломления с соотношением $n_f > n_s > n_c$ Свет, попадая в волновод, испытывает ПВО последовательно от обеих границ и оказывается заключен в волноводной среде и рапространяется вдоль нее с максимальной эффективностью

\begin{figure}[h!]
\includegraphics[width=0.9\textwidth]{img/planar_propagation.pdf}
\caption{Распространение света в волноводе}
\end{figure}
